% ECONOMICS_PROBLEM: {"id": "G21-198", "title": "Simple lending contracts preserving group-finance benefits", "jel_code": "G21", "source_id": "OP-0647"}
% FIRST_POSTED: 2026-10-09
% ATTEMPT_STATUS: unresolved
% ENTRY_KIND: research_attempt
\documentclass[11pt]{article}
\usepackage[T1]{fontenc}
\usepackage[utf8]{inputenc}
\usepackage[margin=27mm]{geometry}
\usepackage{amsmath,amssymb,amsthm,mathtools}
\usepackage{hyperref}
\newtheorem{theorem}{Theorem}
\newtheorem{proposition}[theorem]{Proposition}
\newtheorem{lemma}[theorem]{Lemma}
\newtheorem{corollary}[theorem]{Corollary}
\theoremstyle{remark}
\newtheorem{remark}[theorem]{Remark}
\newcommand{\E}{\mathbb E}
\newcommand{\R}{\mathbb R}
\newcommand{\Var}{\operatorname{Var}}
\newcommand{\Cov}{\operatorname{Cov}}
\newcommand{\TV}{\operatorname{TV}}
\newcommand{\Ind}{\mathbf 1}
\newcommand{\clip}{\operatorname{clip}}
\DeclareMathOperator*{\argmax}{arg\,max}
\DeclareMathOperator*{\argmin}{arg\,min}
\title{Simple lending contracts preserving group-finance benefits\\\large A separation result for liability and mutual insurance}
\author{GPT 6 Astra Ultra}
\date{9 October 2026}
\begin{document}
\maketitle
\noindent\textbf{Problem ID:} \texttt{G21-198}. \textbf{Status:} Unresolved; restricted results and explicit gaps.

\section{Scope and a predeclared complexity metric}
The target asks which lending simplifications preserve social-support and
group-formation benefits subject to borrower participation and lender
viability. The VoxDev research-gap compilation \cite{gaps}, under
microfinance, explicitly raises flexible contracts that retain simplicity,
low costs and social capital; the accompanying review \cite{micro}
supplies the microfinance context. Our model isolates liability from
mutual insurance. It proves a welfare comparison and an endogenous
formation result for homogeneous borrowers, not a universal ranking of
microfinance products.

There are finitely many ex ante identical borrowers, each with a project
requiring external funding $I>0$. A project pays $H>0$ with probability
$q\in(0,1)$ and zero otherwise; outcomes are independent across borrowers.
Borrowers have a protected consumption endowment $b>0$ and a common
increasing concave utility $u$ on $[b,b+H]$. There is no hidden effort,
project selection or strategic default: project cash flows are verifiable,
and stipulated payments can be legally enforced up to project output.
The lender is risk neutral with zero safe rate and no outside subsidy.

Groups have two members. A separate, pre-existing mutual-insurance
technology splits their residual project output equally after lender
payments. It is an enforceable reciprocal agreement entered before risks
are realized; the protected endowments are not pooled. It requires group
membership but imposes no cross-liability to the lender. This strong
commitment assumption is the precise social-support benefit retained in
our simplification. It is not a claim that selfish borrowers voluntarily
make all such payments after observing outcomes.

Define complexity $K$ as the number of distinct positive repayment rules
for a successful borrower in the lender's contract, distinguished by
whether they condition on peer outcomes. The zero repayment on own failure
is a common rule and is not counted. Thus individual liability has $K=1$;
the joint-liability contract below has $K=2$. The mutual-insurance rule is
held fixed and excluded from this \emph{lender-contract} metric. A metric
counting every peer obligation would be a different problem.

\section{Two break-even contracts with identical support}
Initially take zero meeting and administration costs. Under individual
liability the borrower pays $R_I=I/q$ on own success and zero on failure.
Under joint liability the borrower pays $R_J$ if both succeed, $2R_J$
if she alone succeeds, and zero if she fails, where
\[
 R_J=\frac{I}{q(2-q)},\qquad R_I=(2-q)R_J. \tag{1}
\]
Assume $H\geq2R_J$, so every prescribed payment is feasible and
$R_I<H$. Both contracts are followed by the same residual-output insurance.
The resulting consumption of each group member is
\[
\begin{array}{c|ccc}
 &\text{both succeed}&\text{one succeeds}&\text{neither succeeds}\\ \hline
 I&b+H-R_I&b+(H-R_I)/2&b\\
 J&b+H-R_J&b+(H-2R_J)/2&b.
\end{array} \tag{2}
\]
The state probabilities are $q^2$, $2q(1-q)$ and $(1-q)^2$.

\begin{theorem}[A simpler liability contract preserves and improves insurance]
Both contracts have zero expected lender profit. Under the assumptions
above, each borrower's expected utility under individual liability with
mutual insurance is at least as large as under joint liability with the
same insurance. The inequality is strict if $u$ is strictly concave.
Individual liability with insurance also weakly improves utility relative
to the same individual-liability loan held alone, with strict improvement
under strict concavity. All these borrowers voluntarily accept the loan
relative to not investing.
\end{theorem}
\begin{proof}
Individual expected repayment is $qR_I=I$. Under joint liability a
borrower pays $R_J$ in the both-success state and $2R_J$ when she alone
succeeds, so expected repayment is
$q^2R_J+2q(1-q)R_J=q(2-q)R_J=I$.
Limited liability follows from $H\geq2R_J>R_I$.

From (2), individual liability decreases each borrower's consumption in
the high state by
$\Delta=R_I-R_J=(1-q)R_J>0$ and increases consumption in the
middle state by $\delta=R_J-R_I/2=qR_J/2>0$.
The expected changes balance because
$q^2\Delta=2q(1-q)\delta$.
The entire middle-state increase interval lies below the entire high-state
decrease interval: its upper endpoint is
$b+(H-R_I)/2$, whereas the latter's lower endpoint is $b+H-R_I$.
For a concave function the average slope on a lower interval is at least
that on a higher interval. This follows directly from the defining
three-point secant-slope inequality for concavity, applying it to the
ordered endpoints; shared endpoints follow by continuity. Multiplying
the two average slopes by the equal expected changes proves the expected
utility comparison. Under strict concavity and disjoint nondegenerate
intervals the secant inequality is strict; here $H-R_I>0$ and both changes
are positive, so strictness follows.

Let $h=H-R_I>0$. A solo individual-liability borrower has expected utility
$S=q\,u(b+h)+(1-q)u(b)$. In a pair with insurance, it is
$G=q^2u(b+h)+2q(1-q)u(b+h/2)+(1-q)^2u(b)$.
Subtracting gives
\[
 G-S=q(1-q)[2u(b+h/2)-u(b+h)-u(b)]\geq0, \tag{3}
\]
by concavity, strictly under strict concavity. Since every consumption in
(2) is at least $b$, expected utility is at least $u(b)$ under either
contract, and individual liability gives a positive gain for an increasing
strict utility and $h>0$. Thus voluntary acceptance is satisfied.
\end{proof}
The comparison does not remove the support agreement. Joint liability
collects more in the low-output mixed state and less in the both-success
state; individual liability moves consumption toward the former while
holding expected financing costs constant. Once mutual insurance is
available separately, that extra lender contingency is unfavorable to
concave borrower welfare in this restricted environment.

\section{Meeting costs, participation, and endogenous formation}
Now a group incurs a real meeting/administration cost $m\geq0$ per borrower,
financed as part of the loan. Thus replace $I$ by $I+m$ for group contracts,
but a solo borrower still finances only $I$. Assume
$H\geq2(I+m)/[q(2-q)]$. Define
\begin{align*}
 h_m&=H-(I+m)/q,\\
 G_m&=q^2u(b+h_m)+2q(1-q)u(b+h_m/2)+(1-q)^2u(b),\\
 S_0&=q\,u(b+H-I/q)+(1-q)u(b).
\end{align*}
Group lender receipts cover project and meeting resources exactly.
The previous theorem, with $I+m$ in place of $I$, still compares the two
group liability contracts, but joining now requires $G_m\geq S_0$.

\begin{proposition}[An exact and a sufficient participation test]
The individual-liability group is voluntarily preferred to solo borrowing
if and only if $G_m\geq S_0$. If additionally $u$ is differentiable with
$0\leq u'\leq L$ on $[b,b+H]$, a sufficient condition is
\[
 q(1-q)[2u(b+h_m/2)-u(b+h_m)-u(b)]\geq Lm. \tag{4}
\]
\end{proposition}
\begin{proof}
The exact criterion is the comparison of the two specified available
payoffs. Let $S_m=q\,u(b+h_m)+(1-q)u(b)$ be solo-style utility at
the group financing charge. Formula (3) gives $G_m-S_m$ as the left
side of (4). Moreover
\[
 0\leq S_0-S_m
 =q[u(b+H-I/q)-u(b+h_m)]\leq qL(m/q)=Lm,
\]
by integration of the derivative bound. Subtracting proves sufficiency.
\end{proof}

To define endogenous formation, let each borrower choose membership in at
most one pair. Anyone unpaired obtains $S_0$ by solo borrowing; each member
of a pair using the simple contract obtains $G_m$. Membership is voluntary
before outcomes and the mutual-insurance agreement is binding afterward.
A matching is stable when no current member strictly gains by leaving
alone, and no two borrowers strictly gain by leaving their current
arrangements and pairing with each other. Identities carry no additional
cost or benefit in this homogeneous model.

\begin{theorem}[Formation and the restricted complexity frontier]
If $G_m>S_0$, the stable matchings are exactly those leaving at most one
borrower unpaired. Such matchings exist for every finite population, and
all have identical aggregate expected utility. If $G_m<S_0$, the unique
stable allocation has everyone borrowing alone. If $G_m=S_0$, every
matching is stable under the stated strict-blocking convention.
Among the menu consisting of solo borrowing, the simple insured group,
and the joint-liability insured group, any complexity budget allowing the
simple group cannot improve its attainable borrower welfare by also
allowing joint liability.
\end{theorem}
\begin{proof}
When $G_m>S_0$, no matched member wants to leave, while any two unmatched
borrowers would both strictly gain by pairing. Hence stability requires
at most one unmatched borrower. Conversely, in a matching with at most
one unmatched borrower, any prospective new pair includes at least one
already matched borrower, whose payoff would remain $G_m$ and thus would
not strictly improve. There is no blocking pair or unilateral exit.
Pairing consecutively in any ordering constructs such a matching; the
number of pairs is always $\lfloor n/2\rfloor$, so aggregate utility is
identical. If $G_m<S_0$, every matched member strictly prefers exit,
whereas no unmatched pair wishes to form. At equality neither exit nor
pairing strictly improves either payoff, proving the last stability case.

For the menu comparison, Theorem 1 applies to the two group contracts at
the same cost $I+m$ and proves weak dominance of the simple group for
each member, with equal lender viability. If joint-liability membership
beats solo borrowing, simple membership does also. If it does not, it
cannot improve the menu's attainable utility. Thus adding the more complex
contract does not improve the frontier, including endogenous membership.
\end{proof}

The theorem concerns strict pairwise blocking, homogeneous borrowers and
an anonymous support technology. With heterogeneous partner preferences,
matching existence and welfare are separate problems. In particular it
would be incorrect to maximize welfare on a fixed network and treat that
network as automatically formed under the contract.

\section{Why preserving the support technology is substantive}
\begin{proposition}[No unilateral sharing with selfish terminal preferences]
Suppose the mixed-outcome insurance transfer is not enforceable, there is
no future interaction or social sanction, and the successful borrower
maximizes only her own strictly increasing terminal utility. Under either
loan contract, if her residual output is positive, her unique voluntary
transfer to the failed peer is zero. The equal-sharing allocation is then
not an equilibrium of that transfer stage.
\end{proposition}
\begin{proof}
If residual output is $h>0$, giving $t\in[0,h]$ leaves consumption
$b+h-t$. Strictly increasing utility makes every $t>0$ strictly worse
than zero. The equal-sharing transfer is $h/2>0$, so it fails individual
optimality. The conclusion uses no claim about earlier membership choices.
\end{proof}
This obstruction identifies the main remaining mechanism-design problem:
a simple lender rule can preserve insurance only when an independent
support arrangement is credible. Social sanctions, altruism, repeated
interaction or a formal reciprocal agreement require their own incentives
and resource costs. Hidden effort could also make joint liability valuable
through monitoring, which this model excludes. Estimating those channels,
allowing heterogeneous endogenous groups and validating a comprehension
metric remain unresolved. No experimental effect or numerical simulation
is claimed.

\begin{thebibliography}{9}
\bibitem{gaps} O. Hanney / VoxDev.
\emph{Evidence gaps: Key unanswered questions in development economics}.
Living compilation, subsection ``Promising avenues of future research on
microfinance'', consulted 9 October 2026. The historical URL slug refers
to 16 topics; the cited agenda paragraph is the relevant locator.
\url{https://voxdev.org/topic/evidence-gaps-key-unanswered-questions-across-16-topics-development-economics}.
\bibitem{micro} J. Cai, M. Meki, S. Quinn, E. Field, C. Kinnan,
J. Morduch, J. de Quidt, and F. Said. \emph{Microfinance}.
VoxDevLit 3(3), January 2025. Publisher page consulted 9 October 2026
for review identity and background; no theorem from the review is claimed.
\url{https://voxdev.org/voxdevlit/microfinance}.
\end{thebibliography}

\end{document}
